\section{Flag sebagai Dasar Pengambilan Keputusan}
\label{sec:flag-keputusan}

Sampai bab ini, seluruh program yang dibahas dieksekusi secara
berurutan (\eng{sequential execution}): unit kendali menjalankan
instruksi pada alamat \reg{CS}:\reg{IP}, lalu memajukan \reg{IP} ke
instruksi berikutnya. Program yang berguna hampir selalu membutuhkan
lebih dari itu: memilih di antara beberapa tindakan, mengulang
pekerjaan, dan berhenti ketika suatu syarat terpenuhi.

\subsection{Mengapa Prosesor Membutuhkan Flag}

Prosesor 8086 tidak memiliki instruksi ``jika'' seperti pada bahasa
tingkat tinggi, melainkan sekumpulan \term{flag}---bit penanda
keadaan pada \reg{FLAGS}---dan instruksi lompatan yang memeriksanya.
Pengambilan keputusan pun selalu berlangsung dalam dua langkah: satu
instruksi mengubah keadaan \eng{flag} sesuai hasil operasi, lalu satu
instruksi lompatan bersyarat memutuskan apakah alamat eksekusi
dipindahkan.

Langkah kedua tidak menghitung apa pun: bila tidak ada instruksi yang
menyiapkan \eng{flag} sebelumnya, lompatan bersyarat akan memakai
\eng{flag} sisa dari operasi yang jauh di belakang. Register
\reg{FLAGS} beserta makna setiap bitnya telah dibahas pada
Bab~\ref{chap:register}; lima \eng{flag} yang paling sering dipakai
untuk percabangan dirangkum pada Tabel~\ref{tab:flag-keputusan}.

\begin{table}[H]
    \centering
    \caption{Lima flag utama pada pengambilan keputusan.}
    \label{tab:flag-keputusan}
    \small
    \begin{tabular}{@{}p{1.4cm}p{5.0cm}p{7.0cm}@{}}
        \toprule
        \textbf{Flag} & \textbf{Makna} & \textbf{Dipakai oleh} \\
        \midrule
        \flag{ZF} & Bernilai 1 bila hasil operasi sama dengan nol &
        \asm{JE}, \asm{JNE}, \asm{JG}, \asm{JL}, \asm{LOOPE} \\
        \addlinespace
        \flag{CF} & Bernilai 1 bila terjadi pinjaman atau limpahan
        pada operasi tak bertanda &
        \asm{JB}, \asm{JAE}, \asm{JC}, \asm{JNC} \\
        \addlinespace
        \flag{SF} & Menyalin bit paling kiri (bit tanda) hasil operasi &
        \asm{JS}, \asm{JNS}, serta lompatan bertanda \\
        \addlinespace
        \flag{OF} & Bernilai 1 bila hasil melampaui jangkauan bilangan
        bertanda &
        \asm{JO}, \asm{JNO}, \asm{JG}, \asm{JL} \\
        \addlinespace
        \flag{PF} & Menyatakan kegenapan banyak bit 1 pada byte hasil &
        \asm{JP}, \asm{JNP} \\
        \bottomrule
    \end{tabular}
\end{table}

\important{Lompatan bersyarat tidak membandingkan dua bilangan. Ia
hanya membaca \eng{flag}. Karena itu, instruksi penyiap \eng{flag}
harus ditulis tepat sebelum lompatan.}

\subsection{Instruksi CMP: Menyiapkan Flag melalui Pengurangan}

Instruksi \asm{CMP} (\eng{compare}) melakukan pengurangan antara
operand pertama dan operand kedua, tetapi \emph{tidak menyimpan}
hasilnya; yang tersimpan hanyalah perubahan pada \eng{flag}. Bentuk
dan aturan operandnya sama persis dengan \asm{SUB} yang dibahas pada
Bab~\ref{chap:aritmatika}.

\begin{lstlisting}[language=MASM, numbers=none]
MOV AX, 0005h       ; AX = 5
MOV BX, 0003h       ; BX = 3
CMP AX, BX          ; hitung 5 - 3 = 2, hasil dibuang
; keadaan flag: ZF = 0, SF = 0, CF = 0, OF = 0
JG  LEBIH_BESAR     ; lompat karena 5 > 3 (bertanda)
\end{lstlisting}

Setelah \lstasm|CMP AX, BX| dijalankan, isi \reg{AX} dan \reg{BX} tidak
berubah; yang berubah hanya \reg{FLAGS}. Inilah sebabnya \asm{CMP}
disebut instruksi \term{penyiap flag}. Bila kedua operand sama,
\flag{ZF} bernilai 1; bila operand pertama lebih kecil secara tak
bertanda, terjadi pinjaman sehingga \flag{CF} bernilai 1. Perilaku
inilah yang menjadi dasar seluruh tabel lompatan pada
Subbab~\ref{sec:jcc-loop}.

\subsection{Instruksi TEST: Menyiapkan Flag melalui AND}

Instruksi \asm{TEST} melakukan operasi \asm{AND} bit demi bit, lalu
membuang hasilnya dan hanya menyimpan \eng{flag}. Berbeda dari
\asm{CMP}, \asm{TEST} tidak melakukan pengurangan; ia
\emph{menyaring} bit.

\begin{lstlisting}[language=MASM, numbers=none]
TEST AL, 01h        ; uji bit ke-0 dari AL
JZ   GENAP          ; ZF = 1 berarti bit ke-0 nol, bilangan genap

TEST AL, AL         ; uji apakah AL sama dengan nol
JZ   NOL
\end{lstlisting}

Pemakaian \lstasm|TEST AL, AL| adalah idiom yang sering muncul:
operand dibandingkan dengan dirinya sendiri sehingga \flag{ZF} dan
\flag{SF} mengikuti nilai operand, tanpa mengubah isi \reg{AL}.
Perbedaan kedua instruksi penyiap \eng{flag} ini dirangkum pada
Tabel~\ref{tab:cmp-test}.

\begin{table}[H]
    \centering
    \caption{Perbedaan \asm{CMP} dan \asm{TEST}.}
    \label{tab:cmp-test}
    \small
    \begin{tabular}{@{}p{3.0cm}p{5.6cm}p{5.8cm}@{}}
        \toprule
        \textbf{Aspek} & \textbf{\asm{CMP}} & \textbf{\asm{TEST}} \\
        \midrule
        Operasi dasar & Pengurangan
        (\eng{dst} $-$ \eng{src}) & \asm{AND} bit demi bit
        (\eng{dst} \& \eng{src}) \\
        \addlinespace
        Hasil disimpan & Tidak; hanya \eng{flag} &
        Tidak; hanya \eng{flag} \\
        \addlinespace
        Flag utama yang disiapkan & \flag{ZF}, \flag{CF},
        \flag{SF}, \flag{OF}, \flag{AF}, \flag{PF} &
        \flag{ZF}, \flag{SF}, \flag{PF}; \flag{CF} dan \flag{OF}
        selalu 0 \\
        \addlinespace
        Kegunaan khas & Membandingkan dua nilai, menentukan urutan &
        Menguji bit, memeriksa nol atau tanda \\
        \addlinespace
        Contoh & \code{CMP AX, BX} & \code{TEST AL, 80h} \\
        \bottomrule
    \end{tabular}
\end{table}

\begin{peringatanbox}
Memakai \asm{TEST} untuk membandingkan dua bilangan adalah kekeliruan
yang sering terjadi. \lstasm|TEST AX, BX| tidak memberi tahu apakah
\reg{AX} lebih besar dari \reg{BX}; ia hanya memberi tahu apakah
\emph{ada} bit yang bernilai 1 pada kedua operand secara bersamaan.
Untuk membandingkan nilai, gunakan \asm{CMP}.
\end{peringatanbox}

\subsection{Lompatan Tak Bersyarat JMP dan Jarak Lompatan}

Instruksi \asm{JMP} memindahkan kendali tanpa syarat, misalnya
\lstasm|JMP HITUNG|. Prosesor 8086 tidak menyimpan alamat mutlak
tujuan, melainkan \term{perpindahan} (\eng{displacement}) relatif
terhadap instruksi berikutnya; besar perpindahan inilah yang
menentukan \term{jarak lompatan}.

\begin{enumerate}
    \item \term{Short jump} menyimpan perpindahan dalam satu byte
          bertanda, sehingga jangkauannya hanya \dec{127} byte ke
          depan dan \dec{128} byte ke belakang.

    \item \term{Near jump} menyimpan perpindahan dalam 16 bit,
          sehingga tujuan dapat berada di mana saja dalam segmen kode
          yang sama. Instruksi ini berukuran tiga byte.

    \item \term{Far jump} menyatakan alamat tujuan secara lengkap
          sebagai pasangan \reg{CS}:\reg{IP}, sehingga tujuan dapat
          berada pada segmen kode yang berbeda. Instruksi ini
          berukuran lima byte.
\end{enumerate}

Ketiga jangkauan tersebut dilukiskan pada Gambar~\ref{fig:jarak-lompatan}.

\begin{figure}[H]
    \centering
    \begin{tikzpicture}[font=\small]
        % bingkai segmen kode
        \draw[draw=obeGray, line width=0.7pt, fill=obeGrayLight]
            (0,0) rectangle (2.6,-7.0);
        \node[font=\scriptsize\bfseries, text=obeGray] at (1.3,-0.35)
            {segmen kode};

        % pita jangkauan short
        \fill[obeTealLight] (0,-3.0) rectangle (2.6,-4.0);
        \draw[draw=obeTeal, line width=0.6pt] (0,-3.0) rectangle (2.6,-4.0);

        % titik asal
        \fill[obeMaroon] (2.6,-3.5) circle (2.4pt);

        % keterangan jangkauan short
        \node[anchor=east, font=\scriptsize, text=obeTeal, align=right]
            at (-0.2,-3.5) {jangkauan\\\textbf{short}\\$\pm$ 127 byte};

        % kotak tujuan
        \node[draw=obeTeal, fill=obeTealLight, rounded corners=2pt,
              font=\scriptsize, align=center, inner sep=4pt]
            (ts) at (5.6,-3.5) {tujuan \textbf{short}\\(1 byte perpindahan)};
        \node[draw=obeBlue, fill=obeBlueLight, rounded corners=2pt,
              font=\scriptsize, align=center, inner sep=4pt]
            (tn) at (5.6,-6.1) {tujuan \textbf{near}\\(16 bit, segmen sama)};
        \node[draw=obeAmber, fill=obeAmberLight, rounded corners=2pt,
              font=\scriptsize, align=center, inner sep=4pt]
            (tf) at (5.6,-0.9) {tujuan \textbf{far}\\(segmen kode lain)};

        \draw[-{Stealth[length=2.2mm]}, obeTeal, line width=0.9pt]
            (2.6,-3.5) -- (ts.west);
        \draw[-{Stealth[length=2.2mm]}, obeBlue, line width=0.9pt]
            (2.6,-3.5) -- (tn.west);
        \draw[-{Stealth[length=2.2mm]}, obeAmber, line width=0.9pt]
            (2.6,-3.5) to[out=60, in=180] (tf.west);

        \node[anchor=west, font=\scriptsize, text=obeMaroon]
            at (2.75,-3.5) {\asm{JMP} di sini};
    \end{tikzpicture}
    \caption{Tiga jenis jarak lompatan pada instruksi \asm{JMP} dan
    jangkauan masing-masing.}
    \label{fig:jarak-lompatan}
\end{figure}

Perbedaan ketiganya memengaruhi jumlah byte yang dipakai, seperti
ditunjukkan pada Tabel~\ref{tab:jarak-lompatan}; ukuran ini penting
ketika program harus muat dalam memori terbatas.

\begin{table}[H]
    \centering
    \caption{Jenis jarak lompatan dan pengaruhnya pada ukuran kode.}
    \label{tab:jarak-lompatan}
    \small
    \begin{tabular}{@{}p{1.7cm}p{3.6cm}p{3.0cm}p{2.4cm}p{2.6cm}@{}}
        \toprule
        \textbf{Jenis} & \textbf{Jangkauan} &
        \textbf{Perpindahan} & \textbf{Ukuran} & \textbf{Penulisan} \\
        \midrule
        \eng{short} & $\pm$ 127 byte dari instruksi berikutnya &
        1 byte bertanda & 2 byte & \code{JMP SHORT L1} \\
        \addlinespace
        \eng{near} & Seluruh segmen kode (64 KB) &
        16 bit & 3 byte & \code{JMP L1} \\
        \addlinespace
        \eng{far} & Segmen kode mana pun &
        \reg{CS}:\reg{IP} lengkap & 5 byte &
        \code{JMP FAR PTR L1} \\
        \bottomrule
    \end{tabular}
\end{table}

Assembler memilih jarak lompatan secara otomatis; operator
\code{SHORT} atau \code{FAR PTR} memaksa bentuk tertentu. Bila
programmer memaksa \code{SHORT} padahal jaraknya terlalu jauh,
assembler melaporkan kesalahan \file{Relative jump out of range}.

\begin{infobox}
Pada program besar---misalnya rutin BIOS yang harus muat dalam
64 KB---pemilihan jarak lompatan dapat menghemat ribuan byte. Inilah
contoh nyata bahwa pemrograman tingkat rendah memberi kendali yang
tidak tersedia pada bahasa tingkat tinggi \cite{bartlett2004}.
\end{infobox}
